\SourceLeaf{343}

Der Mittelpunkt des Umkreises

\[
= \frac{2\alpha\beta\gamma(\alpha+\beta+\gamma)}{(\beta+\gamma)(\gamma+\alpha)(\alpha+\beta)}.
\]

Produkt aus dem Radius und dem Betrag von $(\beta+\gamma)(\gamma+\alpha)(\alpha+\beta)$

\[
= 2.
\]

Die Clisen von $\dfrac{\beta\gamma}{\alpha}$, $\dfrac{\alpha\gamma}{\beta}$, $\dfrac{\alpha\beta}{\gamma}$ entsprechen den drei Radien dieses Kreises.

\medskip
\noindent 3. Vier Zahlen $a$, $b$, $c$, $d$ sind einander harmonisch zugeordnet, wenn das
Produkt der Summen der beiden Paare doppelt so groß ist wie die Summe der
Produkte der beiden Paare, oder wenn

\[
(a+b)(c+d) = 2ab+2cd.
\]

Sind die Zahlen reell, so stellen sie vier Punkte auf einer Geraden dar, die
demnach harmonisch geteilt ist.

Mit dieser Gleichung ist verbunden

\[
(c-a)(d-b)+(d-a)(c-b) = 0.
\]

\Needspace{9\baselineskip}Jede Zahl lässt sich aus den drei übrigen mit den folgenden Formeln berechnen

\[
\begin{aligned}
a &= \frac{bc+bd-2cd}{-d-c+2b} = \frac{b(c+d)-2cd}{2b-(c+d)} = \frac{2bb-2cd}{2b-(c+d)}-b \\
  &= +b-\frac{2(c-b)(d-b)}{2b-c-d} \\
  &= +c-\frac{(d-c)(c-b)}{2b-c-d}.
\end{aligned}
\]

Sind $a$, $b$, $c$, $d$ komplexe Zahlen, so liegen die von ihnen dargestellten Punkte
auf einem Kreis, und die Produkte der Seiten $ac.bd$ und $ad.bc$ sind gleich.
